\subsection{拓扑群的覆叠}

命题 4. --- 设 (X, e) 是霍普夫空间（IV, 第~373~页，定义 2），m: X × X → X 是其合成法则。假设空间 X 连通且局部道路连通。设 X' 是 X 的连通覆叠；以 p 表示其投影，并令 e' 为纤维 \(X'_{e}\) 中的一点。

\begin{enumerate}
\def\labelenumi{\alph{enumi})}
\item
  覆叠 X' 是伽罗瓦覆叠，且其自同构群是交换的。
\item
  存在唯一的连续合成法则 \(m' \colon X' \times X' \to X'\)，使得 \(p \circ m' = m \circ (p, p)\) 且 \(m'(e', e') = e'\)。赋予带基点空间 \((X', e')\) 以此合成法则后，它是一个霍普夫空间。
\item
  以由 \(m'\) 诱导的合成法则赋予纤维 \(X'_e\) 后，它是以 \(e'\) 为单位元的群。从 \(\pi_1(X, x)\) 到 \(X'_e\) 的映射 \(\gamma \mapsto e' \cdot \gamma\) 是满射群同态，其核为 \(p_*(\pi_1(X', e'))\)。映射 \(g \mapsto g(e')\) 是从群 \(\text{Aut}_X(X')\) 到群 \(X'_e\) 的同构。
\item
  群 \(\pi_{1}(X,e)\) 是交换的（命题 3）。由于 X 的覆叠 \(X'\) 连通，它于是是伽罗瓦覆叠，且群 \(\operatorname{Aut}_{\mathrm{X}}(\mathrm{X}')\) 同构于商群 \(\pi_{1}(\mathrm{X},e)/p_{*}(\pi_{1}(\mathrm{X}',e'))\)（III, 第~312~页，命题 2 的推论 4）；此群是交换的。
\item
  记 \(q\colon X^{\prime}\times X^{\prime}\to X\) 为映射 \(m\circ (p,p)\)。所需的映射 \(m'\) 是 \(q\) 到 \(X^{\prime}\) 的连续提升，使得 \(m'(e',e') = e'\)。空间 \(X^{\prime}\times X^{\prime}\) 局部道路连通；因此，为证明存在这样的连续提升，只需验证 \(q_{*}(\pi_{1}(X^{\prime}\times X^{\prime},(e^{\prime},e^{\prime})))\) 包含于 \(p_{*}(\pi_{1}(X^{\prime},e^{\prime}))\)（III, 第~308~页，命题 1）。根据命题 3，映射 \(\pi_1(m,(e,e))\) 可识别为群 \(\pi_1(X,e)\) 的合成法则，条件是将 \(\pi_1(X\times X,(e,e))\) 识别为 \(\pi_1(X,e)\times \pi_1(X,e)\)。因此，群 \(q_{*}(\pi_{1}(X^{\prime}\times X^{\prime},(e^{\prime},e^{\prime})))\) 等于群 \(p_{*}(\pi_{1}(X^{\prime},e^{\prime}))\)，从而存在 \(m'\)。由于 \(X^{\prime}\times X^{\prime}\) 连通，\(m'\) 的唯一性由 I, 第~34~页的命题 11 的推论 1 得出。
\end{enumerate}

证明映射 \(m'\) 赋予带基点空间 \((\mathrm{X}', e')\) 霍普夫空间结构。令 \(m_1 \colon X \to X\) 为映射 \(g \mapsto m(g, e)\)，并令 \(\sigma_1 \colon X \times I \to X\) 为在 x 处的带基点同伦，将 \(m_1\) 与 \(Id_X\) 连接起来。置 \(\tau = \sigma_1 \circ (p, \mathrm{Id_I}) \colon X' \times I \to X\)。映射 \(m'_1 \colon X' \to X'\) 由 \(h \mapsto m'(h, e')\) 定义，并提升映射 \(\tau(\cdot, 0)\)；令 \(\tau' \colon X' \times I \to X'\) 为 \(\tau\) 的提升，其初始映射为 \(m'_1\)（III, 第~301~页，命题 2）。映射 \(t \mapsto \tau'(e', t)\) 是到 \(X'\) 的常值映射 \(t \mapsto e\) 的提升；由于 \(\tau'(e', 0) = e'\)，故有 \(\tau'(e', t) = e'\) 对所有 \(t \in I\) 都成立。映射 \(\tau'(\cdot, 1)\) 随后是映射 p 到 \(X'\) 的提升，它将 \(e'\) 映到 \(e'\)。由于 \(X'\) 连通，\(Id_{X'}\) 是从 \(X'\) 到自身且固定点 \(e'\) 的唯一 X-态射；因此 \(\tau'(\cdot,1)=\mathrm{Id}_{\mathrm{X}'}\)。这表明 \(\tau'\) 是在 \(e'\) 处的带基点同伦，将映射 \(m_{1}'\) 与映射 \(Id_{X'}\) 连接起来。同样，存在在 \(e'\) 处的带基点同伦，将映射 \(m_{2}'\colon h\mapsto m'(e',h)\) 与映射 \(Id_{X'}\) 连接起来。这证明了带基点空间 ( \(X',e'\) ) 赋予合成法则 \(m'\) 后是霍普夫空间。

\begin{enumerate}
\def\labelenumi{\alph{enumi})}
\setcounter{enumi}{2}
\tightlist
\item
  由于 \(X'\) 是 X 的连通覆叠，\(e'\) 的轨道映射由 \(\pi_{1}(\mathrm{X},e)\) 在 \(X'_{e}\) 上的作用诱导，是满射，并诱导出从 \(\pi_{1}(\mathrm{X},e)/p_{*}(\pi_{1}(\mathrm{X}',e'))\) 到 \(X'_{e}\) 的双射（III, 第~305~页，定理 1）。
\end{enumerate}

设 \(c\) 和 \(d\) 是 X 中在 \(e\) 处的圈，\(\gamma\) 和 \(\delta \in \pi_1(\mathrm{X},e)\) 是它们的类。设 \(c'\) 和 \(d'\) 是以 \(e'\) 为起点、位于 \(\mathrm{X}'\) 上且分别提升 \(c\) 和 \(d\) 的道路。有 \(e' \cdot \gamma = c'(1)\) 且 \(e' \cdot \delta = d'(1)\)（III, 第~304~页）。根据 IV, 第~374~页的命题 3，圈 \(m \circ (c,d)\) 与圈 \(c*d\) 严格同伦；因此有 \(e' \cdot (\gamma\delta) = e' \cdot (m \circ (c,d))\)。现在，道路 \(m' \circ (c',d')\) 是起点为 \(e'\) 的道路 \(m \circ (c,d)\) 的提升；于是有

\[
e ^ {\prime} \cdot (\gamma \delta) = m ^ {\prime} (c ^ {\prime} (1), d ^ {\prime} (1)) = m ^ {\prime} (e ^ {\prime} \cdot \gamma , e ^ {\prime} \cdot \delta).
\]

因此，映射 \(\gamma \mapsto e' \cdot \gamma\) 是从 \(\pi_1(\mathrm{X}, e)\) 到集合 \(\mathrm{X}'_e\)（其合成法则由 \(m'\) 诱导）的同态。于是，\(\mathrm{X}'_e\) 在由 \(m'\) 诱导的合成法则下是一个群，且 \(e'\) 的轨道映射是从商群 \(\pi_1(\mathrm{X}, e)/p_*(\pi_1(\mathrm{X}', e'))\) 到 \(\mathrm{X}'_e\) 的同构。

断言 c) 的最后一部分随即由 III, 第~311~页的推论 3 得出。

命题 5. --- 保持命题 4 的记号与假设。

\begin{enumerate}
\def\labelenumi{\alph{enumi})}
\item
  如果 m 是结合的（分别为交换的）合成法则，那么 \(m'\) 也是如此。
\item
  如果 e 是作用 m 的右单位元（分别为左单位元），那么 e' 是作用 m' 的右单位元（分别为左单位元）。
\item
  如果 X 是拓扑群，m 是其合成法则，e 是其单位元，则合成法则 m' 赋予 X' 以与 X' 的拓扑相容的群结构，且 e' 是其单位元。映射 p: X' → X 是群同态，其核是离散的并包含于 X' 的中心。
\item
  假设法则 m 是结合的。那么，从 \(X' \times X' \times X'\) 到 X 的、分别把 \((h_{1}, h_{2}, h_{3})\) 映为 \(m(p(h_{1}),m(p(h_{2}),p(h_{3})))\) 和 \(m(m(p(h_{1}),p(h_{2})),p(h_{3}))\) 的映射相等。映射 \((h_{1},h_{2},h_{3})\mapsto m'(h_{1},m'(h_{2},h_{3}))\) 和 \((h_{1},h_{2},h_{3})\mapsto m'(m'(h_{1},h_{2}),h_{3})\) 从 \(X'\times X'\times X'\) 到 \(X'\) 是连续提升，并在点 \((e',e',e')\) 处重合。由于 \(X'\times X'\times X'\) 连通，它们相等（I, 第~34~页，命题 11 的推论 1），这说明法则 \(m'\) 是结合的。
\end{enumerate}

现在假设法则 \(m\) 是交换的；映射 \((h_1, h_2) \mapsto m'(h_1, h_2)\) 和 \((h_1, h_2) \mapsto m'(h_2, h_1)\) 是到 \(X'\) 的连续提升，提升映射 \((h_1, h_2) \mapsto m(p(h_1), p(h_2))\)，并在点 \((e', e')\) 处重合。由于 \(X' \times X'\) 连通，它们相等（同上），且法则 \(m'\) 是交换的。

如果 \(e\) 是法则 \(m\) 的右单位元（分别为左单位元），则映射 \(h \mapsto m'(h, e')\)（分别为 \(h \mapsto m'(e', h)\)）是从 \(X'\) 到 \(X'\) 的覆叠的 X-态射，与 \(\text{Id}_{X'}\) 在点 \(e'\) 处重合；由于 \(X'\) 连通，它因而在 \(X'\) 的每一点处都与之重合（同上），从而得到 \(b\)。

最后证明 c)。设 X 是拓扑群。根据上文，法则 \(m'\) 是结合的，且 \(e'\) 是单位元。令 \(i: X \to X\) 为映射 \(g \mapsto g^{-1}\)。它连续（TG, III, 第~1~页），且同态 \(\pi_{1}(i,e): \pi_{1}(X,e) \to \pi_{1}(X,e)\) 恰好是映射 \(\gamma \mapsto \gamma^{-1}\)（IV, 第~374~页，命题 3）。于是，子群 \((i \circ p)_{*}(\pi_{1}(X',e'))\) 等于 \(p_{*}(\pi_{1}(X',e'))\)。根据 III, 第~308~页的命题 1，存在连续映射 \(i': X' \to X'\)，使得 \(p \circ i' = i \circ p\) 且 \(i'(e') = e'\)。映射 \(h \mapsto m'(h,i'(h))\) 和 \(h \mapsto m'(i'(h),h)\) 从 \(X'\) 到 \(X'\) 是到 \(X'\) 的提升，它们提升从 \(X'\) 到 X 的值为 e 的常值映射。因此它们是常值的，且其像为 \(e' = m'(e',e')\)。于是，\(X'\) 的每个元素 h 都可逆，逆元为 \(i'(h)\)，这表明赋予 \(X'\) 以合成法则 \(m'\) 后它是一个群。根据法则 \(m'\) 的构造，映射 p: \(X' \to X\) 是群同态。由于映射 \(m'\) 和 \(i'\) 连续，\(X'\) 的群结构与其拓扑相容（TG, III, 第~1~页）。纤维 \(p^{-1}(e)\) 是 \(X'\) 的离散子群，包含于 \(X'\) 的中心（I, 第~100~页，推论 3），且 X 同构于商拓扑群 \(X'/p^{-1}(e)\)。

推论。--- 设 G 是连通且局部道路连通的拓扑群。设 G' 是 G 的连通覆叠，p 是其投影，且 e' 是 G 的单位元 e 的纤维 N 中的一个元素。给 G' 赋予唯一的连续合成法则 m': G' × G' → G'，使得 \(p \circ m' = m \circ (p, p)\) 且 \(m'(e', e') = e'\)。如果 i: N → G' 表示规范嵌入，那么 (G', p, i) 是 G 以 N 为核的中心扩张（A, I, 第~63~页）。

